\documentclass[10pt,a4paper]{amsart}
\usepackage[margin=19mm]{geometry}
\usepackage{amsmath,amssymb,mathtools}
\usepackage[scr=rsfs,cal=euler]{mathalfa}
\usepackage{xfrac}
\usepackage[unicode,hidelinks,bookmarks=false]{hyperref}
\newcommand{\ie}{i.e.\ }
\newcommand{\cf}{cf.\ }
\newcommand{\resp}{respectively\ }
\newcommand{\Subsection}{\subsection}
\providecommand{\nobreakdash}{\nobreak\hbox{-}}
\setlength{\emergencystretch}{2em}
\multlinegap0pt
\usepackage{amssymb}
\newtheorem{assumption}{Assumption}
\newtheorem{proposition}{Proposition}
\newtheorem{lemma}{Lemma}
\def \E{\mathbf{E}}
\def \P{\mathbf{P}}
\title{Mean-Field Backward Stochastic Differential Equations with Nonlinear Resistance and Double Mean Reflections}
\author{Hanwu Li, Jin Shi}
\date{Source: arXiv:2605.15781v1 (CC BY 4.0)}
\begin{document}
\maketitle

\noindent\emph{Source and licence.} Mean-Field Backward Stochastic Differential Equations with Nonlinear Resistance and Double Mean Reflections. Version arXiv:2605.15781v1 (CC BY 4.0), distributed by arXiv under the Creative Commons Attribution 4.0 International licence, http://creativecommons.org/licenses/by/4.0/, as recorded in the arXiv licence metadata for this exact version and shown on the abstract page https://arxiv.org/abs/2605.15781v1. Full licence notice: \texttt{licenses/arxiv-2605.15781-CC-BY-4.0.txt}.\par\smallskip

\noindent\textit{Editorial scope.} Counted unit: the article's lemma le3-2 (source0.tex lines 547--549). The article's complete original proof of it is delivered with it (source0.tex lines 551--621, one \texttt{proof} environment of 71 lines). No mathematics is written, repaired or re-derived: the statement and the proof are the article's own lines, in the article's own notation, and no retained line is substituted or re-flowed.  Retained uncounted as the source-local context the proof needs: lines 336--339; lines 458--469; lines 471--479; lines 482--484; lines 498--508; lines 532--539.  Nothing was omitted from any retained span, so the package carries no omission record.  No retained line contains a citation, so the wrapper carries no bibliography and no bibliography entry.  No retained line contains a figure environment, an image-inclusion command or drawing code, so no image is delivered and the package declares no asset dependency.  Editorial formatting only, in addition: an AMS \texttt{amsart} preamble that restates the article's own declarations and macros used by the retained lines, namely the environments \texttt{assumption}, \texttt{proposition}, \texttt{lemma} on separate counters, together with the standard packages those lines need.  The retained environments print in this document as Assumption 1, Proposition 1, Lemma 1 in order of appearance, whereas the article prints the same items with the numbering of its own full document.  The complete native source of the article as distributed in the arXiv e-print for this exact version is bundled unchanged as \texttt{sources/200682/source0.tex}.
\begin{equation}\label{diffK}
\sup_{t\in[0,T]}\left|K^1_t-K^2_t\right|
\leq  2\frac{C}{c}\left|a^1-a^2\right|+4\frac{C}{c}\sup_{t\in[0,T]}\left|s^1_t-s^2_t\right|+\frac{2}{c}(\bar{L}_T\vee\bar{R}_T),
\end{equation}

\begin{assumption}\label{assLR}
\begin{itemize}
\item[(1)] For any fixed $(\omega,x)\in \Omega\times\mathbb{R}$, $L(\omega, \cdot, x), R(\omega, \cdot, x)$ are continuous;%for any fixed $x\in\mathbb{R}$, $L(\cdot,\cdot,x), R(\cdot,\cdot,x)$ is progressively measurable and  
\item[(2)]  There exists a constant $M>0$ such that $$\mathbf{E}\left[\sup_{t\in[0,T]}\left|L(t,0)\right|+\sup_{t\in[0,T]}\left|R(t,0)\right|\right]\leq M.$$
\item[(3)] For any fixed $(\omega,t)\in \Omega\times [0,T]$, $L(\omega,t,\cdot),R(\omega,t,\cdot)$ are strictly increasing and there  exist two constants $c,C$ satisfying $0<c<C$ such that for any $x,y\in \mathbb{R}$,
\begin{align*}
&c\left|x-y\right|\leq \left|L(\omega,t,x)-L(\omega,t,y)\right|\leq C\left|x-y\right|,\\
&c\left|x-y\right|\leq \left|R(\omega,t,x)-R(\omega,t,y)\right|\leq C\left|x-y\right|;
\end{align*}
\item[(4)] $\inf_{\omega,t,x} \left(R(\omega,t,x)-L(\omega,t,x)\right)>0$.
\end{itemize}
\end{assumption}

\begin{assumption}\label{ass1ipterminalp}
\begin{itemize}
      \item[(i)]  The terminal value $\xi \in \mathcal{L}^{p}_T $ satisfies $\mathbf{E}[L(T,\xi)]\leq 0 \leq \mathbf{E}[R(T,\xi)]$; 
      \item[(ii)] The generator $f$: $[0,T]\times\Omega\times \mathbb{R}\times \mathcal{P}_1(\mathbb{R})\times\mathbb{R}\times \mathcal{P}_1(\mathbb{R})\times \mathbb{R}$ $\rightarrow$ $\mathbb{R}$ is a $\mathcal{P}\times\mathcal{B}(\mathcal{P}_1(\mathbb{R}))\times\mathcal{B}(\mathbb{R})\times\mathcal{B}(\mathcal{P}_1(\mathbb{R}))\times\mathcal{B}(\mathbb{R})$-measurable map, and there exists a constant $\lambda>0$ such that for all $t\in[0,T]$, $y, y', k, k' \in \mathbb{R}$, $\mu, \mu' \in \mathcal{P}_1(\mathbb{R})$, $z, z'\in \mathbb{R}$, $\nu, \nu'\in\mathcal{P}_1(\mathbb{R})$, 
\begin{align*}
\left|f(t,y,\mu,z,\nu,k)-f(t,y',\mu',z',\nu',k')\right|\leq \lambda\left(\left|y-y'\right|+\left|z-z'\right|+W_1(\mu,\mu')+W_1(\nu,\nu')+\left|k-k'\right|\right).
\end{align*}
   \end{itemize}
\end{assumption}

\begin{assumption}\label{assgk}
For any $t\in[0,T]$, given $y\in BV[0,T]$,  we have $G_{t}(y)=G_{t}\left ( \left \{ y_{s\wedge t} \right \}_{0\le s\le t\le T}  \right ) $ and  $G_t(0)=0$. Moreover, there exists a constant $\lambda>0$, such that for any $y,y'\in BV[0,T]$,   $$\left | G_{t}(y)-G_t(y') \right | \leq \lambda \sup _{u\in[0,t]}\left | y_u-y'_{u} \right |.$$
\end{assumption}

\begin{proposition}\label{pro-1}
  Let  Assumption \ref{ass1ipterminalp} and \ref{assLR} hold. Given $C\in \mathcal{H}^p$, the BSDE with double mean reflections 
   \begin{equation}\label{MFBSDEC}
   \begin{cases}
Y_t=\xi+\int_t^T C_s ds-\int_t^T Z_s dB_s+K_T-K_t, \\
\E[L(t,Y_t)]\leq 0\leq \E[R(t,Y_t)], \\
K_t=K^{R}_t-K^{L}_t, \int_0^T \E[R(t,Y_t)]dK_t^{R}=\int_0^T \E[L(t,Y_t)]dK^{L}_t=0,
\end{cases}
\end{equation}
has a unique solution $\left(Y,Z,K\right)\in \mathcal{S}^p\times\mathcal{H}^{p}\times BV[0,T]$.
\end{proposition}

\begin{equation}\label{yzk}
\begin{cases}
Y_t^{y,z,k}=\xi+\int_t^T f(s,y_s,\P_{y_s},z_s,\P_{z_s}, G_s(k))ds-\int_t^T Z_s^{y,z,k} dB_s+K_T^{y,z,k}-K_t^{y,z,k}, \\
\E[L(t,Y_t^{y,z,k})]\leq 0\leq \E[R(t,Y_t^{y,z,k})], \\
K^{y,z,k}_t=K^{y,z,k,R}_t-K^{y,z,k,L}_t, \\
\int_0^T \E[R(t,Y_t^{y,z,k})]dK_t^{y,z,k,R}=\int_0^T \E[L(t,Y_t^{y,z,k})]dK^{y,z,k,L}_t=0
\end{cases}
\end{equation}

\begin{lemma}\label{le3-2}
Let $\varepsilon$ satisfy $\beta(\lambda,c,C,p)\max\left(\varepsilon^{\frac{p}{2}},\varepsilon^p\right)<1$, where $\beta$ is a positive constant depending only on $\lambda, C,c,p$. For $T\in (0,\varepsilon]$, the solution map $\Gamma$ is contractive. 
\end{lemma}

\begin{proof}
In this proof,  $\beta$ always represents a positive constant depending on $\lambda,C,c,p$, which may vary from line to line. Let $(Y^i,Z^i,K^i):=(Y^{y^i,z^i,k^i},Z^{y^i,z^i,k^i},K^{y^i,z^i,k^i})$ be the solution to BSDE  with double mean reflections \eqref{yzk} associated with data $(y^i,z^i,k^i)$, $i=1,2$.

Set
\begin{align*}
&\hat{F}_t=F^1_t-F^2_t, \textrm{ where } F=Y,Z,K,y,z,k, \\ &\hat{f}_t=f(t,y^1_t,\P_{y^1_t},z^1_t,\P_{z^1_t},G_t(k^1))-f(t,y^2_t,\P_{y^2_t},z^2_t,\P_{z^2_t},G_t(k^2)).
\end{align*}
It is easy to check that 
\begin{align*}
|\hat{Y}_t|&=\left|\E_t\left[\int_t^T \hat{f}_sds\right]+\hat{K}_T-\hat{K}_t\right|\\
&\leq \lambda \E_t\left[\int_0^T \left(|\hat{y}_s|+|\hat{z}_s|+\E[|\hat{y}_s|]+\E[|\hat{z}_s|]+\sup_{r\in[0,s]}|\hat{k}_r|\right)ds \right]+|\hat{K}_T|+|\hat{K}_t|.
\end{align*}
Applying the Doob's maximal inequality, we obtain that  
\begin{equation}\begin{split}\label{differY}
 \E\left[\sup_{t\in[0,T]}|\hat{Y}_t|^p\right]
 &\leq \beta \E\left[\left(\int_{0}^{T}|\hat{y}_s|+|\hat{z}_s|+\E[|\hat{y}_s|]+\E[|\hat{z}_s|]+\sup_{r\in[0,s]}|\hat{k}_r|ds\right)^p\right]+\beta \sup_{t\in[0,T]}|\hat{K}_t|^p.\\   
\end{split}
\end{equation}
Recalling the proof of Proposition \ref{pro-1} and \eqref{diffK}, we have
\begin{align*}
\sup_{t\in[0,T]}|\hat{K}_t|\leq  \beta \left\{\sup_{t\in[0,T]}|\hat{s}_t|+\sup_{(t,x)\in[0,T]\times\mathbb{R}}|\hat{l}(t,x)|\vee \sup_{(t,x)\in[0,T]\times\mathbb{R}} |\hat{r}(t,x)|\right\},
\end{align*}
where $\hat{s}_t=s^1_t-s^2_t$, $\hat{l}(t,x)=l^1(t,x)-l^2(t,x)$, $\hat{r}(t,x)=r^1(t,x)-r^2(t,x)$ and for $i=1,2$, 
\begin{align*}
&s^i_t=\E\left[\int_0^t f\left(s,y^i_s,\P_{y^i_s},z^i_s,\P_{z^i_s},G_s(k^i)\right)ds\right], \\ 
&l^i(t,x)=\E\left[L\left(t,\widetilde{Y}^i_t-\E[\widetilde{Y}^i_t]+x\right)\right],\\  &r^i(t,x)=\E\left[R\left(t,\widetilde{Y}^i_t-\E[\widetilde{Y}^i_t]+x\right)\right].
\end{align*}
Here, $\widetilde{Y}^i$ is the first component of the solution to the $\text{BSDE}$ with terminal value $\xi$ and constant driver $f^i$, $i=1,2$, where $$f^i_s=f\left(s,y_s^i,\P_{y_s^i},z_s^i,\P_{z^i_s},G_s(k^i)\right).$$ By Assumptions \ref{assLR} and \ref{assgk}, noting that $\widetilde{Y}^i_t=\E_t\left[\xi+\int_t^T f\left(s,y_s^i,\P_{y_s^i},z_s^i,\P_{z^i_s},G_s(k^i)\right)d s\right]$, we obtain that 
\begin{align*}
\left|l^1(t,x)-l^2(t,x)\right|&\leq C\E\left[\left|\left(\widetilde{Y}^1_t-\E[\widetilde{Y}^1_t]\right)-\left(\widetilde{Y}^2_t-\E[\widetilde{Y}^2_t]\right)\right|\right]\leq 2C\E[|\widetilde{Y}^1_t-\widetilde{Y}^2_t|]\\
&\leq \beta \E\left[\int_0^T |\hat{y}_s|+|\hat{z}_s|+\E[|\hat{y}_s|]+\E[|\hat{z}_s|]+\sup_{r\in[0,s]}|\hat{k}_r|d s\right ].
\end{align*}
Similar estimates hold for $|s^1_t-s^2_t|$ and $|r^1(t,x)-r^2(t,x)|$. 
Applying H{\"o}lder inequality, the above analysis yields that 
\begin{align}\label{differK}
\sup_{t\in[0,T]}|\hat{K}_t|^p\leq \beta\E\left[\left(\int_0^T |\hat{y}_s|+|\hat{z}_s|+\E[|\hat{y}_s|]+\E[|\hat{z}_s|]+\sup_{r\in[0,s]}|\hat{k}_r|ds\right)^p \right].
\end{align}
Combining Eqs. \eqref{differY} and \eqref{differK} yields that 
\begin{align*}
\E\left[\sup_{t\in[0,T]}|\hat{Y}_t|^p\right]\leq \beta\E\left[\left(\int_0^T |\hat{y}_s|+|\hat{z}_s|+\E[|\hat{y}_s|]+\E[|\hat{z}_s|]+\sup_{r\in[0,s]}|\hat{k}_r|d s \right)^p\right].
\end{align*}
Note that we have
\begin{align*}
\int_0^t \hat{Z}_s dB_s=\int_0^t\hat{f}_s ds+\hat{Y}_t-\hat{Y}_0+\hat{K}_t.
\end{align*}
Simple calculation yields that
\begin{align*}
\E\left[\sup_{t\in[0,T]}\left|\int_0^t \hat{Z}_sd B_s\right|^p\right]&\leq \beta\E\left[\sup_{t\in[0,T]}\left|\int_0^t\hat{f}_s ds\right|^p+\sup_{t\in[0,T]}|\hat{Y}_t|^p+\sup_{t\in[0,T]}|\hat{K}_t|^p\right]\\
&\leq \beta\E\left[\left(\int_0^T |\hat{y}_s|+|\hat{z}_s|+\E[|\hat{y}_s|]+\E[|\hat{z}_s|]+\sup_{r\in[0,s]}|\hat{k}_r| ds\right)^p\right],
\end{align*}
Combining with the B-D-G inequality yields that 
\begin{align*}
   \E\left[\left(\int_0^T |\hat{Z}_s|^2 d s\right)^{\frac{p}{2}}\right]&\leq \beta\E\left[\left(\int_0^T |\hat{y}_s|+|\hat{z}_s |+\E[|\hat{y}_s|]+\E[|\hat{z}_s|]+\sup_{r\in[0,s]}|\hat{k}_r| ds\right)^p\right].
\end{align*}
By the H\"{o}lder inequality and the Fubini theorem, we finally deduce that 
\begin{align*}
&\E\left[\sup_{s\in[0,T]}|\hat{Y}_s|^p+\left(\int_0^T |\hat{Z}_s|^2 d s\right)^{\frac{p}{2}}+\sup_{t\in[0,T]}|\hat{K}_t|^p\right]\\
\leq &\beta\E\left[\left(\int_0^T \left(|\hat{y}_s|+|\hat{z}_s|+\E[|\hat{y}_s|]+\E[|\hat{z}_s|]+\sup_{r\in[0,s]}|\hat{k}_r|\right)ds\right)^p\right]\\
\leq &\beta\E\left[\left(\int_0^T |\hat{y}_s|+|\hat{z}_s|+\E[|\hat{y}_s|]+\E[|\hat{z}_s|]ds\right)^p+\left(\int_0^T\sup_{r\in[0,s]}|\hat{k}_r|ds\right)^p\right]\\
\leq &\beta\max\left(T^{\frac{p}{2}},T^p\right)\E\left[\sup_{t\in[0,T]}|\hat{y}_t|^p+\left(\int_0^T|\hat{z}_s|^2 ds\right)^{\frac{p}{2}}\right]+ T^p \sup_{r\in[0,T]}|\hat{k}_r|^p\\
\leq &\beta\max\left(T^{\frac{p}{2}},T^p\right)\E\left[\sup_{t\in[0,T]}|\hat{y}_t|^p+\left(\int_0^T|\hat{z}_t|^2 dt\right)^{\frac{p}{2}}+\left(\sup_{t\in[0,T]}|\hat{k}_t|\right)^p\right].
\end{align*}
Therefore, we can choose an appropriate  $T \in(0,\varepsilon]$ such that
$$
\begin{aligned}
  &\left\|Y^{1}-Y^{2}\right\|_{\mathcal{S}^{p}}+\left\|Z^{1}-Z^{2}\right\|_{\mathcal{H}^{p}}+\sup_{t\in[0,T]}\left|K_{t}^{1}-K_{t}^{2}\right|\\
  &\leq \beta\max\left(T^{\frac{p}{2}},T^p\right) \left(\left\|y^1-y^2\right\|_{\mathcal{S}^{p}}+\left\|z^1-z^2\right\|_{\mathcal{H}^{p}}+\sup_{t\in[0,T]}\left|k_t^1-k_t^2\right|\right). %\beta\max\left(T^{\frac{p}{2}},T^p\right)
\end{aligned}
$$
The proof is complete.
\end{proof}

\end{document}
